\chapter{Experimental Design}
\label{chap:design}

The purpose of the experiment is to compare the complete behaviour of two
tests, not just their p-values on convenient tables. Each design point fixes
two population distributions in advance. Repeated independent samples from
those populations show how often each procedure returns a valid result and
how often it rejects. Because the population MI values are known, a rejection
at equal MI is a false positive; a rejection at unequal MI is detection.

\section{A population pair before the design grid}

Consider two \(2\times2\) populations whose row and column margins are
\((0.7,0.3)\) for \(P\) and \((0.8,0.2)\) for \(Q\). The construction produces,
to the displayed precision,
\begin{equation}
 P\simeq
 \begin{pmatrix}0.53273&0.16727\\0.16727&0.13273\end{pmatrix},
 \qquad
 Q\simeq
 \begin{pmatrix}0.68175&0.11825\\0.11825&0.08175\end{pmatrix}.
 \label{eq:design-example-pair}
\end{equation}
These are \emph{probability} tables, not observed counts. Their target MI
values are 0.02 and 0.03 nats, so the true magnitude of the difference is
0.01 nats. Sampling new count tables changes the estimated MI values, but
does not move this fixed population pair to a new horizontal-axis location.
Calculations use the stored unrounded probabilities; the rounded display in
Equation~\eqref{eq:design-example-pair}, determine the samples.

\section{Constructing fixed joint distributions}

For one population, let \(a=(a_1,\ldots,a_r)\) and
\(b=(b_1,\ldots,b_c)\) be prescribed row and column probabilities. The
independence table is \(B=ab^{\mathsf T}\), with \(B_{ij}=a_i b_j\).
Association is introduced by
\begin{equation}
 R(t)=ab^{\mathsf T}+tH,
 \qquad H\mathbf{1}_c=0,
 \qquad \mathbf{1}_r^{\mathsf T}H=0.
 \label{eq:additive-construction}
\end{equation}
The zero row and column sums mean that changing \(t\) leaves the margins
fixed. For the main construction, \(H\) is zero except in the first two rows
and columns, where it is
\begin{equation}
 H_{1:2,1:2}=\begin{pmatrix}1&-1\\-1&1\end{pmatrix}.
 \label{eq:first-change-block}
\end{equation}
Thus \(t\) is added to the two diagonal cells and removed from the other two
cells. The changes cancel within each affected row and column. The same
matrix specifies a \(2\times2\) test and a localised change in a larger
table; only its embedding differs. This follows the minimal four-violation
sparse alternative used by \citet[Section~3(b)]{berrett2021}. Here it is adapted to compare two fixed
populations: their margins are prescribed separately and the perturbation
strength is chosen to attain a specified MI and is not used directly as the
horizontal-axis effect.

The largest nonnegative strength before a cell with \(H_{ij}<0\) reaches zero
is
\begin{equation}
 t_{\max}=\min_{H_{ij}<0}\frac{a_i b_j}{-H_{ij}}.
 \label{eq:design-tmax}
\end{equation}
The construction solves for \(0\le t<t_{\max}\) such that \(I\{R(t)\}\)
equals the requested target. Targets at or beyond the endpoint are not
admitted. This endpoint is specific to the chosen path, not the largest
possible MI over every table with those margins. Both constructed tables are
checked for strictly positive cells, correct margins and target MI before any
sampling. Appendix~\ref{app:construction} gives the numerical details.

For a dimension of size \(k\), the skewed margin used in the design is
\begin{equation}
 m_k(\rho)=\left(\rho,\frac{1-\rho}{k-1},\ldots,
                   \frac{1-\rho}{k-1}\right).
 \label{eq:skew-margin}
\end{equation}
The three principal profiles are specified separately for the two
populations in Table~\ref{tab:margin-profiles}. For a rectangular
\(r\times c\) table, the formula is applied with \(k=r\) to row margins and
\(k=c\) to column margins. At a zero MI difference, the same-margin profiles
produce \(P=Q\) under this construction; the different-skew profile has
\(I(P)=I(Q)\) but \(P\ne Q\).

\begin{table}[htbp]
\centering
\caption{Row and column margins used for each population in the principal
design. ``Uniform'' means the vector with every entry equal to \(1/k\).}
\label{tab:margin-profiles}
\begin{tabular}{@{}lll@{}}
\toprule
Profile & \(P\): rows and columns & \(Q\): rows and columns\\
\midrule
Uniform & uniform & uniform\\
Same skew & \(m_k(0.8)\) & \(m_k(0.8)\)\\
Different skew & \(m_k(0.7)\) & \(m_k(0.8)\)\\
\bottomrule
\end{tabular}
\end{table}

In the usual effect direction, \(I(P)=I_0\) and
\(I(Q)=I_0+|\Delta|\). The paired solve chooses separate strengths for
\(P\) and \(Q\); it does not add random noise to their true MI values. The
reverse-direction imbalance experiment assigns the higher target to \(P\).
The statistical convention remains \(\Delta=I(P)-I(Q)\), which is negative in
the usual direction. Figures use the nonnegative magnitude \(|\Delta|\).

Two focused checks change the pattern or constructor while keeping the
interpretation of target MI\@. A ``rare block'' puts Equation
\eqref{eq:first-change-block} in the final two rows and columns. A ``spread''
pattern is the outer product of equally spaced, centred row and column scores;
its row and column sums are zero. The ordinal log-linear comparator fits the
prescribed margins while varying an ordinal interaction strength to meet the
target MI\@. Matching margins and MI across constructors does not imply matching
all cell probabilities or expected counts.

\section{Sampling, tests and measured outcomes}

For each unique configuration, 20,000 independent pairs of multinomial count
tables are drawn from its fixed \(P\) and \(Q\). Normal Wald and Expanded
Welch are applied to the \emph{same} count-table pair. This pairing makes
method differences more precisely measurable than unrelated simulations,
though it does not remove Monte Carlo uncertainty. The two-sided nominal
threshold is \(\alpha=0.05\) throughout the statistical study.

Let \(A_m=1\) if method \(m\) returns a valid p-value no larger than 0.05,
and \(A_m=0\) otherwise. The primary plotted rejection rate is
\begin{equation}
 \widehat r_m=\frac{1}{20{,}000}\sum_{s=1}^{20{,}000} A_{m,s}.
 \label{eq:unconditional-rate}
\end{equation}
At \(|\Delta|=0\) this estimates the false-positive rate. At
\(|\Delta|>0\) it estimates the probability of detecting the specified
difference at the fixed nominal threshold, often called power. A rate near
0.05 is desirable under the null; a smaller rate is not automatically better.
Under an alternative, higher detection is useful only alongside acceptable
null calibration and a usable test result.

An invalid calculation contributes no rejection to Equation
\eqref{eq:unconditional-rate}. This convention reports the behaviour of the
complete procedure, but an invalid result is not a valid acceptance of the
null. The companion valid-result rate is the fraction of replicates with a
usable p-value. The saved data also give rejection conditional on each
method's valid outputs and on the subset where both are valid. Those
denominators answer different questions: conditioning each method on its own
valid samples can change which tables are compared. For a single rate, saved
Wilson 95\% intervals \citep{wilson1927} represent Monte Carlo uncertainty at the fixed
population pair. Paired method differences use the same replicates and their
saved paired standard errors. These intervals do not describe uncertainty
over all possible population distributions.

The ``minimum expected count'' diagnostic is
\(\min_{ij}n_Pp_{ij}\) for \(P\), or \(\min_{ij}n_Qq_{ij}\) for \(Q\).
It is not the expectation of the smallest observed count. No lower bound on
that diagnostic was used to admit statistical configurations. Empty observed
cells and rows remain possible. In the plots a hollow marker warns that fewer
than 90\% of outputs of that method were valid; 90\% is a display convention,
not a validated safety threshold.

\section{Experiment families and computation}

Table~\ref{tab:design-families} records the question answered by each family.
The common population construction above applies unless the named factor
changes. Exact parameter sets are in Appendix~\ref{app:construction}, and the
machine-readable protocol is identified in Appendix~\ref{app:reproducibility}.
There are 534 distinct population pairs \((P,Q)\), each specifying two
fixed joint probability tables. Reusing a pair at different sample sizes
creates distinct sampling configurations. There are 3,111 unique statistical
configurations but 4,001 display points:
some configurations are deliberately shown in more than one comparison.
Each unique configuration is sampled once, giving 62,220,000 table pairs in
the final statistical run. The within-family counts in
Table~\ref{tab:design-families} do not sum to 3,111 because the same
configuration can occur in several families. Display reuse is not an
independent replication.

\begin{table}[htbp]
\centering
\caption{The fixed experiment families. All statistical points use 20,000
paired replicates and the 5\% threshold.}
\label{tab:design-families}
\begin{tabularx}{\textwidth}{@{}lrrX@{}}
\toprule
Family & Displays & Unique & Factor examined\\
\midrule
Main & 648 & 648 & Square table size, three margin profiles and equal sample size.\\
Baseline & 648 & 648 & Starting positive MI of 0.0001, 0.001 or 0.02 nats.\\
Changed cells & 144 & 144 & First, rare or spread association patterns.\\
Imbalance & 1,152 & 1,056 & Sample allocation and which population has higher MI.\\
Larger effects & 162 & 162 & MI differences up to 0.2 nats where feasible.\\
Rectangular & 108 & 108 & \(2\times3\) and \(3\times5\) tables.\\
Construction & 432 & 432 & Additive versus ordinal log-linear population tables.\\
Extreme skew & 315 & 315 & Dominant margins up to 0.9995, including \(n=1\) stress.\\
Rare-cell stress & 140 & 140 & First- versus rare-block changes, including \(n=1\) stress.\\
Independence & 144 & 144 & Exact \(I(P)=0\) boundary.\\
Convergence & 108 & 108 & Fixed equal-MI populations with sample sizes up to 50,000.\\
Runtime & 60 & 60 & Complete-call timing regimes, separate from simulation.\\
\bottomrule
\end{tabularx}
\end{table}

The runtime experiment measures a full method call after count tables are
already available. It varies table shape, profile, sample size and null or
alternative status. Each regime uses 200 saved inputs, 20 warm-up calls and
five timed passes. Thus it measures method computation, not the cost of
constructing populations or generating random tables. Median call times and
interquartile ranges are reported in Chapter~\ref{chap:results}. The finite
tested range includes samples of size one in the stress check and an
\(8\times8\) largest square table; it does not imply valid inference for all
sparse distributions or unbounded alphabet sizes.
